\chapter{\IfLanguageName{dutch}{Implementatie}{Implementation}}%

\label{ch:implementatie}

Dit hoofdstuk beschrijft hoe RustProv stapsgewijs werd opgebouwd tijdens de verschillende sprints van het project. Waar het hoofdstuk Methodologie vooral het plan van aanpak op hoofdlijnen toelicht, ligt de focus hier op de concrete technische uitwerking. Per sprint wordt aangegeven welke functionaliteiten werden gerealiseerd, welke ontwerpkeuzes daarbij belangrijk waren en wat de tool op dat moment reeds kon uitvoeren.

\section{Sprint 1: opstarten van virtuele machines}

De eerste implementatiesprint richtte zich op het automatisch kunnen aanmaken en opstarten van een virtuele machine in VirtualBox. Deze fase vormde de technische basis voor de rest van het project, omdat provisioning pas mogelijk wordt zodra een VM correct geregistreerd, geconfigureerd en gestart kan worden. In deze sprint werden de eerste functies \texttt{start} en \texttt{state} voorzien, al waren deze in het begin nog beperkt uitgewerkt.

Tijdens deze sprint werd gebruikgemaakt van \texttt{VBoxManage}-commando's om nieuwe VM's te registreren, de hardware-instellingen te configureren en een bestaand \texttt{.vdi}-bestand als virtuele schijf te koppelen. Daarbij werd een bestaande box gekopieerd naar een tijdelijke doelmap, waarna de UUID van de schijf aangepast werd. Deze stap was nodig om conflicten te vermijden wanneer meerdere VM's op basis van dezelfde box moesten worden aangemaakt. Daarnaast werd ook een shared folder gekoppeld, zodat latere provisioningscripts vanuit de hostomgeving beschikbaar konden worden gemaakt in de gastmachine.

Aan het einde van deze sprint kon RustProv reeds een VM opstarten op basis van een voorbereide box en de toestand van die VM opvragen. De codebasis bestond aanvankelijk uit één groter bestand, maar werd tijdens deze fase al gedeeltelijk herschikt naar een meer modulaire structuur om latere uitbreidingen mogelijk te maken.

\begin{lstlisting}[language=bash,caption={Voorbeeld van console-uitvoer na sprint 1.},label={lst:sprint-1}]
PS C:\Users\JensG\Desktop\gitbphistory\Sprint1_startvm\target\release> .\app.exe start
> VBoxManage createvm --name windows-vm --ostype Windows10_64 --register
Virtual machine 'windows-vm' is created and registered.
UUID: 5e241e88-056f-4145-bd73-52ac2fd87cab
Settings file: 'C:\Users\JensG\VirtualBox VMs\windows-vm\windows-vm.vbox'
> VBoxManage modifyvm windows-vm --memory 2048 --cpus 2 --vram 128 --nic1 nat --audio none
Warning: --audio is deprecated and will be removed soon. Use --audio-driver instead!
> VBoxManage storagectl windows-vm --name SATA Controller --add sata --controller IntelAHCI
> VBoxManage internalcommands sethduuid C:\Users\JensG\.test_dest\windows-vm\vdi.vdi
UUID changed to: 808aab78-851a-4792-aa38-1069a30859e1
> VBoxManage storageattach windows-vm --storagectl SATA Controller --port 0 --device 0 --type hdd --medium C:\Users\JensG\.test_dest\windows-vm\vdi.vdi
> VBoxManage sharedfolder add windows-vm --name shared --hostpath ./scripts --automount
> VBoxManage controlvm windows-vm natpf1 guestssh,tcp,,2222,,22
\end{lstlisting}

De console-uitvoer in Listing~\ref{lst:sprint-1} toont een voorbeeld van het aanmaken en configureren van een Windows-VM met RustProv. De console-uitvoer toont hoe RustProv een Windows-VM aanmaakt en configureert.

\section{Sprint 2: provisioning via SSH}

De tweede sprint richtte zich op het effectieve provisioningproces. Nadat het opstarten van een VM functioneerde, werd de volgende stap het automatisch verbinden met de machine en het uitvoeren van scripts. Hierbij lag de nadruk op netwerktoegang via NAT-portforwarding en op de uitvoering van provisioningsscripts via SSH.

Omdat Linux- en Windowsomgevingen op een andere manier omgaan met scriptuitvoering, werd de provisioningslogica opgesplitst per type gastbesturingssysteem. Voor Linux werd een script uitgevoerd via \texttt{bash} en \texttt{sudo} vanuit de gemounte shared folder. Voor Windows werd een afwijkende aanpak gebruikt, waarbij via SSH een PowerShell-script gestart werd met behulp van \texttt{psexec}. Deze opsplitsing was nodig om binnen één tool toch beide platformen op een consistente manier te kunnen ondersteunen.

Een bijkomend aandachtspunt in deze sprint was het beheer van portforwardingregels. Voor het opzetten van een SSH-verbinding werd een tijdelijke NAT-regel toegevoegd, zodat de host via \texttt{127.0.0.1:2222} de VM kon bereiken. Vooraf werd gecontroleerd of er reeds conflicterende regels bestonden, zodat poortbotsingen vermeden konden worden. Aan het einde van de provisioning werd deze tijdelijke netwerkregel opnieuw verwijderd.

Na deze sprint kon RustProv niet alleen een VM opstarten, maar ook automatisch verbinding maken via SSH en een provisioningscript uitvoeren op zowel Linux- als Windowsmachines. Daarmee werd de kern van de tool functioneel.

\begin{lstlisting}[language=bash,caption={Voorbeeld van console-uitvoer bij Linux-provisioning na sprint 2.},label={lst:sprint-2-linux}]
PS C:\Users\JensG\Desktop\gitbphistory\Sprint2_Linux\target\release> .\app.exe provision
IP address not set for VM 'ubuntu-vm'
VM running on IP: NoneAttempting to provision
Hello, World!
Hello, World from a file in linux!
Provisioning complete! Check the output above for details.
\end{lstlisting}

Listing~\ref{lst:sprint-2-linux} toont het uitvoeren van een provisioningscript op een Linux-machine.

\begin{lstlisting}[language=bash,caption={Voorbeeld van console-uitvoer bij Windows-provisioning na sprint 2.},label={lst:sprint-2-windows}]
PS C:\Users\JensG\Desktop\gitbphistory\Sprint2_windows\target\release> .\app.exe provision
Attempting to provision
Hello, World!
Hello, World from a file in windows!
Provisioning complete! Check the output above for details.
\end{lstlisting}

Listing~\ref{lst:sprint-2-windows} toont het uitvoeren van een provisioningscript op een Windows-machine.

\section{Sprint 3: externe configuratie}

In de derde sprint werd ondersteuning toegevoegd voor configuratie via een extern TOML-bestand. In de eerdere versies van RustProv werden verschillende instellingen nog hardgecodeerd in de broncode. Dat beperkte de flexibiliteit van de tool en maakte het moeilijker om verschillende testscenario's op een reproduceerbare manier te beheren.

Door de introductie van TOML-configuratie konden VM-namen, besturingssysteemtypes, geheugeninstellingen, CPU-toewijzingen, boxkeuzes en bijkomende opties buiten de code worden beschreven. Deze configuratie werd ingelezen en omgezet naar interne datastructuren, waardoor éénzelfde uitvoerbare applicatie verschillende opstellingen kon behandelen zonder hercompilatie. In dezelfde sprint werd ook een \texttt{init}-functie toegevoegd om automatisch de vereiste mappenstructuur voor het project aan te maken.

Na deze sprint kon RustProv op basis van een declaratief configuratiebestand een VM-opstelling interpreteren en de nodige voorbereidingen treffen voor verdere uitvoering. Daarmee werd de tool niet alleen functioneel, maar ook beter bruikbaar voor herhaalde experimenten en verschillende onderwijsopstellingen.

\begin{lstlisting}[language=bash,caption={Voorbeeld van console-uitvoer na het inlezen van TOML-configuratie.},label={lst:sprint-3}]
PS C:\Users\JensG\Desktop\gitbphistory\Sprint2_windows\target\release> .\app.exe start
Starting VM: windowstoml
Image: windows10, Cores: 4, Memory: 4096MB, Scripts: [""]
VM windowstoml does not exist, creating and starting it...
> VBoxManage createvm --name windowstoml --ostype Windows10_64 --register
Virtual machine 'windowstoml' is created and registered.
UUID: 67980b2c-f094-456e-9c4c-7e8bb5dd5749
Settings file: 'C:\Users\JensG\VirtualBox VMs\windowstoml\windowstoml.vbox'
> VBoxManage modifyvm windowstoml --memory 4096 --cpus 4 --vram 128 --nic1 nat --graphicscontroller vmsvga
> VBoxManage storagectl windowstoml --name SATA Controller --add sata --controller IntelAHCI
> VBoxManage internalcommands sethduuid C:\Users\JensG\.test_dest\windowstoml\vdi.vdi
UUID changed to: 106a322d-0d02-4980-866a-50da0ad10a44
> VBoxManage storageattach windowstoml --storagectl SATA Controller --port 0 --device 0 --type hdd --medium C:\Users\JensG\.test_dest\windowstoml\vdi.vdi
> VBoxManage sharedfolder add windowstoml --name shared --hostpath ./scripts --automount
> VBoxManage modifyvm windowstoml --nic2 bridged --bridgeadapter2 Realtek Gaming 2.5GbE Family Controller
> VBoxManage startvm windowstoml --type headless
Waiting for VM "windowstoml" to power on...
VM "windowstoml" has been successfully started.
All VMs started successfully!
\end{lstlisting}

Listing~\ref{lst:sprint-3} toont de console-uitvoer na het inlezen van de
TOML-configuratie. De configuratie bevat informatie over het besturingssysteem,
het aantal CPU-cores, het toegewezen geheugen en de gekoppelde scripts.

\section{Sprint 4: uitbreidingen en gebruiksgemak}

De vierde sprint stond in het teken van uitbreidingen die de tool praktischer maakten in reële gebruiksscenario's. Een eerste belangrijke toevoeging was ondersteuning voor snapshots. Hiermee kon vóór een provisioningsrun een herstelpunt worden gemaakt, zodat bij fouten sneller teruggekeerd kon worden naar een stabiele toestand zonder de volledige VM opnieuw op te bouwen.

Daarnaast werd de code verder opgesplitst zodat bepaalde acties zowel voor één afzonderlijke VM als voor meerdere VM's konden worden toegepast. Ook de command-line interface werd verbeterd met \texttt{clap}, waardoor commando's en parameters duidelijker gestructureerd konden worden. In dezelfde sprint werd een afzonderlijke functie toegevoegd om de status van een VM op te vragen, zodat reeds actieve machines konden worden herkend en overgeslagen indien nodig. Verder werd optionele ondersteuning voorzien voor een bridge adapter als tweede netwerkinterface.

Op het einde van deze sprint kon RustProv niet alleen configuraties inlezen en provisionen, maar ook scenario's met meerdere VM's beter beheren, snapshots aanmaken en meer bruikbare feedback geven aan de gebruiker via de CLI.

\begin{lstlisting}[language=bash,caption={Voorbeeld van console-uitvoer voor restoratie van een snapshots na sprint 4.},label={lst:sprint-4-restore}]
PS C:\Users\JensG\Desktop\gitbphistory\Sprint4_status_and_snapshots\target\release> .\app.exe restore                     
Checking state of VM: ubuntupsnap
Image: ubuntu, Cores: 4, Memory: 4096MB, Scripts: ["setup.sh"]
> VBoxManage controlvm ubuntupsnap poweroff
0%...10%...20%...30%...40%...50%...60%...70%...80%...90%...100%
> VBoxManage snapshot ubuntupsnap restore snapshot1
Restoring snapshot 'snapshot1' (05c0863f-197b-4e8c-8705-4881ff663cbb)
0%...10%...20%...30%...40%...50%...60%...70%...80%...90%...100%
\end{lstlisting}

Listing~\ref{lst:sprint-4-restore} toont het herstellen van een snapshot van
een Linux-VM. Eerst wordt de VM uitgeschakeld, waarna de geselecteerde snapshot
wordt teruggezet.

\begin{lstlisting}[language=bash,caption={Voorbeeld van console-uitvoer voor statuscontrole na sprint 4.},label={lst:sprint-4-state}]
PS C:\Users\JensG\Desktop\gitbphistory\Sprint4_status_and_snapshots\target\release> .\app.exe state
Checking state of VM: ubuntupsnap
Image: ubuntu, Cores: 4, Memory: 4096MB, Scripts: ["setup.sh"]
VM 'ubuntupsnap' is currently running
\end{lstlisting}

Listing~\ref{lst:sprint-4-state} toont het opvragen van de huidige status van
een virtuele machine.

\section{Sprint 5: foutafhandeling en robuustheid}

De laatste sprint richtte zich op foutafhandeling, prioriteitsbeheer en robuuster gedrag in probleemsituaties. Daarbij werd een prioriteitsmechanisme toegevoegd zodat VM's in een vooraf bepaalde volgorde konden worden opgestart en geprovisioned. Dit was relevant voor scenario's waarin bepaalde machines pas correct konden functioneren nadat een andere reeds actief was. Om die reden werd ook gekozen voor \texttt{IndexMap}, zodat de invoervolgorde uit de configuratie behouden bleef.

Daarnaast werden verschillende functies aangepast om correct om te gaan met reeds actieve VM's. De provisioner controleerde eerst de toestand van een machine en sloeg overbodige opstartstappen over wanneer de VM al draaide. Verder werd een \texttt{kill}-functie toegevoegd als noodoplossing voor situaties waarin VirtualBox niet correct reageerde. Ook werden stop- en remove-functionaliteiten uitgebreid met geforceerde varianten, zodat vastgelopen VM's toch beheersbaar bleven.

Een belangrijk onderdeel van deze sprint was de verdere herwerking van Windows-provisioning. In de oorspronkelijke versie bleek scriptuitvoering onder een gewone gebruiker met administratieve rechten niet altijd voldoende. Daarom werd de aanpak aangepast zodat scripts via PSTools en \texttt{psexec} als SYSTEM-gebruiker konden worden uitgevoerd. Hierdoor werd Windows-provisioning betrouwbaarder voor systeemgerichte configuratiestappen.

Na deze sprint beschikte RustProv over een bruikbare proof of concept implementatie die VM's kon aanmaken, opstarten, provisionen, beheren en in foutscenario's gecontroleerd kon reageren. Daarmee was de basis gelegd voor de daaropvolgende experimenten en benchmarkvergelijkingen. 

\begin{lstlisting}[language=bash,caption={Voorbeeld van console-uitvoer na sprint 5 met meerdere VM's en prioriteit.},label={lst:sprint-5-prioriteit}]
PS C:\Users\JensG\Desktop\gitbphistory\Sprint5_priority_and_error\target\release> .\app.exe start --provision  --priority
Starting VM: ubuntu2
Image: ubuntu, Cores: 4, Memory: 4096MB, Scripts: ["setup.sh"]
> VBoxManage createvm --name ubuntu2 --ostype Linux_64 --register
Virtual machine 'ubuntu2' is created and registered.
UUID: 578bc95d-8138-4d2d-9eb4-0b14bf5d78be
Settings file: 'C:\Users\JensG\VirtualBox VMs\ubuntu2\ubuntu2.vbox'
> VBoxManage modifyvm ubuntu2 --memory 4096 --cpus 4 --vram 128 --nic1 nat --graphicscontroller vmsvga
> VBoxManage storagectl ubuntu2 --name SATA Controller --add sata --controller IntelAHCI
> VBoxManage internalcommands sethduuid C:\Users\JensG\.test_dest\ubuntu2\vdi.vdi
UUID changed to: a2e8ca39-185c-4615-aad7-72893c96b90d
> VBoxManage storageattach ubuntu2 --storagectl SATA Controller --port 0 --device 0 --type hdd --medium C:\Users\JensG\.test_dest\ubuntu2\vdi.vdi
> VBoxManage sharedfolder add ubuntu2 --name shared --hostpath ./scripts --automount
> VBoxManage startvm ubuntu2 --type headless
Waiting for VM "ubuntu2" to power on...
VM "ubuntu2" has been successfully started.
provisioning VM...
Provisioning VM: ubuntu2
Image: ubuntu, Cores: 4, Memory: 4096MB, Scripts: ["setup.sh"]
Provisioning VM: ubuntu2 with script: setup.sh
Attempting to provision
> VBoxManage controlvm ubuntu2 natpf1 guestsshrust,tcp,,2222,,22
Connection established!
Hello, World!
Hello, World from a file!
> VBoxManage controlvm ubuntu2 natpf1 delete guestsshrust
Provisioning complete! Check the output above for details.
VM ubuntu2 started and provisioned successfully!
Starting VM: ubuntu1
Image: ubuntu, Cores: 4, Memory: 4096MB, Scripts: ["setup.sh"]
> VBoxManage createvm --name ubuntu1 --ostype Linux_64 --register
Virtual machine 'ubuntu1' is created and registered.
UUID: 4848da28-99c1-4001-9664-2ea10bd43e06
Settings file: 'C:\Users\JensG\VirtualBox VMs\ubuntu1\ubuntu1.vbox'
> VBoxManage modifyvm ubuntu1 --memory 4096 --cpus 4 --vram 128 --nic1 nat --graphicscontroller vmsvga
> VBoxManage storagectl ubuntu1 --name SATA Controller --add sata --controller IntelAHCI
> VBoxManage internalcommands sethduuid C:\Users\JensG\.test_dest\ubuntu1\vdi.vdi
UUID changed to: 692b9d2b-e94f-425e-8262-0961d8f47390
> VBoxManage storageattach ubuntu1 --storagectl SATA Controller --port 0 --device 0 --type hdd --medium C:\Users\JensG\.test_dest\ubuntu1\vdi.vdi
> VBoxManage sharedfolder add ubuntu1 --name shared --hostpath ./scripts --automount
> VBoxManage startvm ubuntu1 --type headless
Waiting for VM "ubuntu1" to power on...
VM "ubuntu1" has been successfully started.
provisioning VM...
Provisioning VM: ubuntu1
Image: ubuntu, Cores: 4, Memory: 4096MB, Scripts: ["setup.sh"]
Provisioning VM: ubuntu1 with script: setup.sh
Attempting to provision
> VBoxManage controlvm ubuntu1 natpf1 guestsshrust,tcp,,2222,,22
Connection established!
Hello, World!
Hello, World from a file!
> VBoxManage controlvm ubuntu1 natpf1 delete guestsshrust
Provisioning complete! Check the output above for details.
VM ubuntu1 started and provisioned successfully!
\end{lstlisting}

Listing~\ref{lst:sprint-5-prioriteit} toont het opstarten en provisionen van
meerdere VM's op basis van hun prioriteit. VM~\texttt{ubuntu2} wordt eerst
opgestart en geprovisioneerd, waarna VM~\texttt{ubuntu1} wordt verwerkt. Dit
komt doordat VM~\texttt{ubuntu2} een hogere prioriteit heeft.

\begin{lstlisting}[language=bash,caption={Voorbeeld van console-uitvoer bij foutafhandeling.},label={lst:sprint-5-error}]
PS C:\Users\JensG\Desktop\gitbphistory\Sprint5_priority_and_error\target\release> .\app.exe start
thread 'main' (27348) panicked at src\commands\tomlread.rs:33:10:
The TOML formatting is invalid!: TomlError { message: "missing field `image`", raw: Some("[ubuntu1]\r\nimage = \"ubuntu\"\r\ncores = 4\r\nmemory = 4096\r\nscript = [\"setup.sh\"]\r\npriority = 2\r\nbridgeNetwork = \"Realtek Gaming 2.5GbE Family Controller\"\r\n\r\n[ubuntu2]\r\ncores = 4\r\nmemory = 4096\r\nscript = [\"setup.sh\"]\r\npriority = 1\r\nbridgeNetwork = \"Realtek Gaming 2.5GbE Family Controller\"\r\n"), keys: [], span: Some(0..0) }
note: run with `RUST_BACKTRACE=1` environment variable to display a backtrace
\end{lstlisting}

Listing~\ref{lst:sprint-5-error} toont de foutafhandeling bij een onvolledige
TOML-configuratie. In de configuratie ontbreekt het verplichte veld
\texttt{image} voor VM~\texttt{ubuntu2}, waardoor RustProv een foutmelding
geeft en de uitvoering stopt.